% Resume style (preamble only). The agent writes the document body from your
% resume. To use your own style, put a .tex file with the same macros in
% profile/resume-templates/one-page.tex
\documentclass[a4paper,10pt]{article}

\usepackage{latexsym}
\usepackage[empty]{fullpage}
\usepackage{titlesec}
\usepackage[usenames,dvipsnames]{color}
\usepackage{verbatim}
\usepackage{enumitem}
\usepackage{hyperref}
\usepackage{fancyhdr}
\usepackage{multicol}
\usepackage{bibentry}
\usepackage{hanging}
\usepackage[T1]{fontenc}
\usepackage{charter}
\usepackage{soul}
\usepackage{fontawesome5}
\usepackage{academicons}


\pagestyle{fancy}
\fancyhf{} % clear all header and footer fields
\fancyfoot{}
\renewcommand{\headrulewidth}{0pt}
\renewcommand{\footrulewidth}{0pt}

% Adjust margins
\addtolength{\oddsidemargin}{-0.6in}
\addtolength{\evensidemargin}{-0.5in}
\addtolength{\textwidth}{1.2in}
\addtolength{\topmargin}{-.7in}
\addtolength{\textheight}{1.4in}

\urlstyle{rm}

\raggedbottom
% Justified text with hyphenation disabled: TeX packs each line as full as it
% can by adjusting word spacing, and never breaks a word across lines.
\hyphenpenalty=10000
\exhyphenpenalty=10000
\emergencystretch=3em
\setlength{\tabcolsep}{0in}

% Global list tightening
\setlist{nosep, topsep=2pt, partopsep=0pt, itemsep=2pt, parsep=0pt}

% Sections formatting
\newcommand{\sectitlefmt}[1]{\so{\MakeUppercase{#1}}}
\titleformat{\section}{\normalsize\bfseries}{}{0em}{\sectitlefmt}[\color{black}\titlerule]
\titlespacing*{\section}{1pt}{7pt}{3pt}

%-------------------------
% Custom commands
\newcommand\publication[1]{%
    \smallskip\par\hangpara{1.5em}{1}\bibentry{#1}\smallskip
}

\newcommand{\resumeItem}[2]{
  \item\small{
    \textbf{#1}{: #2}
  }
}

\newcommand{\resumeItemWithoutTitle}[1]{
  \item\small{
    {\vspace{-2pt}}
  }
}

\newcommand{\resumeSubheading}[4]{
  \item
    \begin{tabular*}{0.97\textwidth}{l@{\extracolsep{\fill}}r}
      \textbf{#1} & #2 \\
      \textit{#3} & \textit{#4} \\
    \end{tabular*}
}

\newcommand{\expSubheading}[4]{
  \item
  \textbf{#1}, \textit{#3} \hfill \textit{#4}
}


\newcommand{\resumeSubItem}[2]{\resumeItem{#1}{#2}}

\renewcommand{\labelitemi}{\textbullet}

\newcommand{\resumeSubHeadingListStart}{\begin{itemize}[leftmargin=0pt, label={}]}
\newcommand{\resumeSubHeadingListEnd}{\end{itemize}}
\newcommand{\resumeItemListStart}{\begin{itemize}[leftmargin=1.5em]}
\newcommand{\resumeItemListEnd}{\end{itemize}}
